\documentclass[withoutpreface,bwprint]{cumcmthesis}

\usepackage[framemethod=TikZ]{mdframed}
\usepackage{url}
\usepackage{subcaption}
\usepackage{geometry}
\usepackage{ulem}
\usepackage{setspace}
\usepackage{array}
\usepackage{graphicx}
\usepackage{multirow}
\usepackage{hyperref}
\usepackage{amsmath}
\usepackage{longtable}
\usepackage{float}
\usepackage{booktabs}
\graphicspath{{figures/}}

\geometry{a4paper,left=25mm,right=25mm,top=25mm,bottom=25mm}

% ---- 标题字号：一级三号黑体居中；二级四号、三级小四黑体 ----
\ctexset{
  section = {format = \centering\zihao{3}\heiti},
  subsection = {format = \zihao{4}\heiti},
  subsubsection = {format = \zihao{-4}\heiti}
}

% ---- 表题居中（表题置表上方，图题置图下方） ----
\captionsetup[table]{justification=centering}

% ---- 版面压缩：缩减浮动体与公式前后间距，消除过多留白 ----
\setlength{\abovecaptionskip}{4pt}
\setlength{\belowcaptionskip}{2pt}
\setlength{\intextsep}{8pt plus 2pt minus 2pt}
\setlength{\textfloatsep}{10pt plus 2pt minus 2pt}
\setlength{\floatsep}{8pt plus 2pt minus 2pt}
\AtBeginDocument{%
  \setlength{\abovedisplayskip}{5pt plus 1pt minus 2pt}%
  \setlength{\belowdisplayskip}{5pt plus 1pt minus 2pt}%
  \setlength{\abovedisplayshortskip}{3pt plus 1pt}%
  \setlength{\belowdisplayshortskip}{3pt plus 1pt}%
}
\title{微网与外部电网电力调控策略}

\begin{document}

\maketitle

\begin{abstract}
\setlength{\parskip}{0pt}
针对光伏出力与负载需求不匹配、光伏预报动态更新以及外网电价波动条件下的微网购电调控问题，本文以满足负载需求与储能运行约束为前提，建立``日前计划---滚动调整---实际执行''相衔接的调度模型，并从储能价值、预报收益和价格影响三个方面量化分析购电费用的变化。

\textbf{针对问题一}，将全天离散为 144 个十分钟时段，建立包含能量平衡、储能递推及日循环约束的线性规划模型，采用\textbf{两阶段词典序优化}：先最小化购电费用，再在费用锁定条件下最小化充放电总量，以剔除无效能量循环。最优方案全天购电量为 \textbf{59,482.70 kWh}，购电费用为 \textbf{35,126.95 元}，较无储能基准降低 26.90\%；结合对偶分析与动态规划的独立数值验证，初始储电量的边际节省价值约为 0.4754 元/kWh。

\textbf{针对问题二}，在当日负载已知的假设下，以过去四天光伏实际均值构造基础预测，并利用历史误差分位数形成\textbf{因果保守裕度}，保证计划仅依赖决策时刻已获得的信息；实际供电缺口按五倍电价紧急购电，储电量跨日连续传递。2025 年 2---12 月共 334 天的总现金购电费用为 \textbf{13,252,341.09 元}，紧急购电量为 \textbf{55,801.05 kWh}，紧急购电费仅占总费用的 2.72\%。

\textbf{针对问题三}，将整点光伏预报重建为十分钟尺度输入，结合日内三次预报更新与执行反馈滚动修正购电和充放电安排，并按最终生效调整量统一结算。评价期总现金购电费用为 \textbf{13,120,194.06 元}，紧急购电量为 \textbf{42,672.22 kWh}；四时刻更新组合较仅使用 0:00 预报\textbf{节省 443,957.78 元}（降幅 3.27\%），其中引入 12:00 预报的增量收益最大，表明日内新增预报信息具有显著经济价值。

\textbf{针对问题四}，在当日电价已知的条件下，沿用前述两类策略并以波动电价重新求解：日前计划策略与滚动调整策略的评价期总费用分别为 \textbf{13,850,454.64 元}和 \textbf{13,727,033.66 元}，较对应固定电价方案分别上升 4.51\% 和 4.63\%。\textbf{交叉结算分解表明}，价格环境变化分别带来约 74.45 万元和 76.00 万元的成本上升，策略重优化仅分别抵消约 14.64 万元和 15.32 万元，价格环境的影响远大于策略重优化的对冲能力。

通过能量守恒、运行边界、信息可得性及输出一致性检验，并结合终端价值、储能效率、预报误差与价格机制的\textbf{灵敏度分析}，本文明确了不同信息与价格条件下的调度特征，为微网协调储能运行和外网购电提供量化依据。

\keywords{微网购电；储能调度；线性规划；滚动调整；费用分解}
\end{abstract}
\clearpage


\section{问题重述与分析}

\subsection{问题背景}

微网由光伏、储能和外部电网共同供能。光伏出力与负荷需求在时间上不完全匹配，储能因而承担跨时段转移电量的作用：在低价或光伏富余时充电，在高价或供能不足时放电。调度既要逐时段满足负荷，又要遵守容量、功率、效率和禁止售电等约束。

相较于单纯的确定性经济调度，本题的核心困难是信息逐步揭示。问题一给出确定性单日基准；问题二令实际光伏在计划时不可见，并以5倍电价惩罚执行缺口；问题三允许日内更新光伏预报并引入偏差计价；问题四进一步加入波动电价。因此，四问应在同一物理系统和一致会计口径下递进比较。

\subsection{研究任务及递进关系}

题目给定储能容量为12,000 kWh，最大充、放电功率均为5,000 kW，储电量运行区间为1,200---10,800 kWh，初始储电量为6,000 kWh。本文将四问的决策权限概括如下。

\begin{table}[!htbp]
\centering\small
\caption{四个问题的信息结构与输出任务}
\resizebox{\textwidth}{!}{%
\begin{tabular}{|c|c|c|c|c|}
\hline
问题 & 决策时可用信息 & 核心决策 & 主要计价规则 & 规定输出 \\
\hline
一 & 附件1电价、负荷和光伏预测 & 单日购电与充放电 & 计划即执行 & \texttt{result1.xlsx} \\
\hline
二 & 固定电价；当日负荷已知；光伏仅用历史预测 & 逐日计划与实际执行 & 计划量全额计费；缺口按5倍价 & \texttt{result2.xlsx} \\
\hline
三 & 0:00、6:00、12:00、18:00光伏预报 & 初始计划与后续调整 & 下调部分按0.5倍价，上调部分按1.5倍价 & \texttt{result3.xlsx} \\
\hline
四 & 波动电价；其余权限分别同问题二、三 & 波动价格下重新优化 & 分别沿用问题二、三规则 & \texttt{result4-2.xlsx}、\texttt{result4-3.xlsx} \\
\hline
\end{tabular}}
\end{table}

其中，“当日负荷在0:00可得”和“问题四当日电价在0:00可得”均为本文对题面数据使用方式的明确假设，不冒充题面直接给定事实；其影响在模型边界与敏感性分析中另行说明。

\subsection{各问分析}

\textbf{问题一}的电价、负荷和光伏预测均已知，且储能状态把144个时段耦合在一起。目标和约束均为线性，故采用线性规划求全局最优；再以最小吞吐量作为次级目标，从等价最优解中选取物理结构更清晰的方案。本问同时提供两把后续标尺：储能相对无储能的节费幅度，以及日末库存的边际价值。

\textbf{问题二}必须区分计划和执行。计划过少会触发5倍电价的紧急购电，计划过多则形成已付费但未调用的电量，二者代价不对称。本文用严格早于决策日的光伏历史构造点预测和保守折减，以逐日线性规划生成计划；执行时使用真实光伏和当前储电量，剩余缺口才由紧急购电补足。该结构使信息可得性、物理流和费用流能够分别核验。

\textbf{问题三}只在问题二基础上增加预报更新和调整权。每次更新只能影响尚未执行的时段，并须同时权衡预测修正收益与0.5/1.5倍偏差结算成本。不同预报时刻的价值应通过同一模型下的逐步加入实验判断，而不能仅按预报次数评价。

\textbf{问题四}只替换电价序列，其余信息权限和物理约束保持不变。通过“固定策略换价格”和“固定价格换策略”的交叉结算，可把总费用变化分解为价格环境效应与策略重优化效应，避免把两种作用混为一谈。

本文技术路线如图1所示。

\begin{figure}[!htbp]
\centering
\includegraphics[width=\textwidth]{fig_paper_route.png}
\caption{本文技术路线图：从典型日最优计划到波动电价决策（四问同标尺）}
\end{figure}

\section{模型假设}

\subsection{题面给定与直接采用的条件}

（1）一天划分为144个十分钟时段，时段内功率视为常数，所有功率均乘以\(\Delta t=1/6\) h换算为电量。

（2）储能容量为12,000 kWh，最大充、放电功率均为5,000 kW，故单时段充、放电量上限为833.333 kWh；储电量保持在1,200---10,800 kWh。本文将题面“充放电效率90\%”解释为充电和放电单向效率均为0.9，因此往返效率为0.81；另一种“往返效率为90\%”的解释仅作为敏感性口径。

（3）微网不向外网售电；供电缺口以交易时刻电价的5倍紧急购电。问题三中，计划量高于最终生效调整量的部分按0.5倍电价计费，调整量高于计划量的部分按1.5倍电价计费。

\subsection{信息可得性假设}

（4）问题二至问题四每天0:00可获得当日负荷曲线。题面仅明确提供全年实际负荷而未提供负荷预报，因此这是本文为隔离光伏不确定性所作的主假设；若负荷同样不可预知，模型费用会上升，相关结果只作边界对照。

（5）尚未发生的实际光伏不参与计划。问题二只使用决策日前已经实现的光伏历史；问题三的预报严格按0:00、6:00、12:00、18:00的发布时间进入模型，且只调整尚未执行的时段。

（6）问题四主方案假设每天0:00可获得附件4的当日电价序列。题面未给价格预报机制，故将给定价格作为计划输入；使用历史价格预测的情形作为信息边界对照。

（7）2025年1月作为历史信息准备期，储能不充不放，初始6,000 kWh保持到2月1日0:00。正式评价期为2025年2月1日至12月31日，共334天；日末实际储电量连续传递到下一日。

\subsection{建模简化及影响}

（8）同一十分钟时段不同时充电与放电；光伏无法被负荷或储能利用时允许弃光。问题二中普通购电按计划量全额结算，未调用部分不返售、不跨时段转移，也不进入实际能量平衡。

（9）问题三的多次调整按各时段最终生效调整量相对0:00原计划结算一次，中间版本不重复计费。除储能转换损耗外，忽略自放电、容量衰减、循环寿命和额外运维成本。因此，本文费用是给定竞赛合同下的购电成本，而非储能全寿命成本。

\section{符号说明}

下标\(d\)表示日期，\(t\)表示日内十分钟时段，\(\tau\in\{0,6,12,18\}\)表示预报发布时刻。主要符号如下。

\begin{table}[!htbp]
\centering\small
\caption{主要符号及单位}
\resizebox{\textwidth}{!}{%
\begin{tabular}{|c|l|c|}
\hline
符号 & 含义 & 单位 \\
\hline
\(T,\Delta t\) & 每日时段数144；单时段长度\(1/6\) & ---，h \\
\hline
\(p_t,p_{d,t}\) & 固定分时电价；逐日逐时段电价 & 元/kWh \\
\hline
\(L_{d,t},E^L_{d,t}\) & 负荷功率；时段负荷电量 & kW，kWh \\
\hline
\(P^{PV}_{d,t},E^{PV}_{d,t}\) & 实际光伏功率；实际光伏电量 & kW，kWh \\
\hline
\(\bar E^{PV}_{d,t},\widehat E^{PV}_{d,t}\) & 四日均值基础预测；保守修正后的光伏预测 & kWh \\
\hline
\(G_t,G^{pl}_{d,t},G^0_{d,t}\) & 问题一购电量；问题二计划量；问题三0:00原始计划量 & kWh \\
\hline
\(A_{d,t},U_{d,t},H_{d,t}\) & 最终生效调整量；实际调用普通电；紧急购电量 & kWh \\
\hline
\(C_{d,t},D_{d,t},S_{d,t}\) & 充电量、放电量、时段末储电量 & kWh \\
\hline
\(R^{PV}_{d,t},W^G_{d,t}\) & 弃光量；已付费但未调用的计划电量 & kWh \\
\hline
\(\eta_c,\eta_d\) & 充电效率、放电效率，均取0.9 & --- \\
\hline
\(r_{d,t},\mathcal R_d,\kappa_d\) & 相对预测误差、成熟误差集、逐日因果保守裕度 & --- \\
\hline
\(q,m\) & 问题三净需求保守度分位、附加光伏裕度 & --- \\
\hline
\(\lambda,\lambda_T\) & 问题二及后续模型的日末库存估值系数，均取0.4720 & 元/kWh \\
\hline
\(J_{plan},J_{adj},J_{emg}\) & 计划购电费、调整相关费用、紧急购电费 & 元 \\
\hline
\(V_0(s)\) & 日末状态固定时、初始储电量为\(s\)的最小购电费 & 元 \\
\hline
\end{tabular}}
\end{table}

其中\((x)^+=\max(x,0)\)。同一符号在计划层和执行层分别以上标\(pl\)和\(ex\)区分；帽号表示预测量。费用比较必须保持评价期、信息集、执行规则和终端处理一致。

\section{模型建立与求解}

\subsection{问题一模型建立与求解}

\subsubsection{时间离散与电量转换}

将一天划分为 \(T = 144\)个时段，每个时段长度为

\begin{equation}\Delta t = \frac{1}{6}\text{~}h.\end{equation}

附件中的时间标签按区间右端点解释：第 1 个数据格对应
\([0\!:\!00,0\!:\!10)\)，可见标签为 0:10；第 144 个数据格对应
\([23\!:\!50,24\!:\!00)\)，可见标签为 0:00+1。结果文件按数据位置依次对应上述时段。

设第 \(t\)时段负载功率和光伏预测功率分别为 \(L_{t}\)和
\({\widehat{P}}_{t}^{PV}\)，对应电量为

\begin{equation}E_{t}^{L} = L_{t}\Delta t,\ {\widehat{E}}_{t}^{PV} = {\widehat{P}}_{t}^{PV}\Delta t.\end{equation}

以下购电量、充放电量、弃光量及储电量均以 kWh 为单位，电价
\(p_{t}\)的单位为元/kWh。

\subsubsection{日前购电优化模型}

设 \(G_{t}\)为外网购电量，\(C_{t},D_{t}\)
分别为储能充电量和放电量，\(R_{t}\) 为弃光量，\(S_{t}\) 为第
\(t\)时段结束时的储电量。

在给定光伏预测条件下，供电与用电满足

\begin{equation}G_{t} + {\widehat{E}}_{t}^{PV} + D_{t} = E_{t}^{L} + C_{t} + R_{t}.\end{equation}

该式将可利用光伏、外网购电和储能放电分配至负载与储能充电，不能利用的光伏电量由
\(R_{t}\)表示。其中

\begin{equation}G_{t} \geq 0,\ 0 \leq R_{t} \leq {\widehat{E}}_{t}^{PV}.\end{equation}

充、放电效率分别取 \(\eta_{c} = \eta_{d} = 0.9\)，储能状态递推为

\begin{equation}S_{t} = S_{t - 1} + \eta_{c}C_{t} - \frac{D_{t}}{\eta_{d}}.\end{equation}

储能运行边界为

\begin{equation}{1200 \leq S_{t} \leq 10800,
}{0 \leq C_{t},D_{t} \leq 5000\Delta t = \frac{5000}{6}.}\end{equation}

依照问题一的日循环要求及给定初值，有

\begin{equation}S_{0} = S_{144} = 6000.\end{equation}

以全天购电费用最小为目标：

\begin{equation}\min J_{1} = \sum_{t = 1}^{144}p_{t}G_{t}.\end{equation}

目标函数与上述约束均为线性形式，因此采用线性规划进行全天联合求解，使储能在满足运行边界的条件下协调各时段购电。

\subsubsection{两阶段词典序求解}

第一阶段求解最小购电费用，得到

\begin{equation}J_{1}^{*} = 35,126.95\text{~元}.\end{equation}

同一最优费用可能对应多组充放电方案。为减少退化解中的无效能量循环，第二阶段在原有约束和第一阶段费用锁定条件下，求解

\begin{equation}\min\sum_{t = 1}^{144}\left( C_{t} + D_{t} \right).\end{equation}

费用锁定按数值容差实施，优先保证经济目标，再降低储能总吞吐量。主模型采用
HiGHS
线性规划求解器，充放电互斥性通过解后检查验证，必要时采用混合整数模型处理。

最终执行方案的购电费用为

\begin{equation}J_{1,actual} = 35,126.95\text{~元}.\end{equation}

其与第一阶段最优值的差约为
\(3.51 \times 10^{- 5}\)元，属于费用锁定的数值容差范围。

\subsubsection{购电与储能调度结果}

全天购电量为 59,482.70 kWh，购电费用为 35,126.95 元。按题目表 1
的格式，指定时段结果如下。

\begin{table}[!htbp]
\centering
\caption{问题一指定时段购电量及全天购电量和购电费}
\small
\setlength{\tabcolsep}{4pt}
\begin{tabular}{|c|c|c|c|c|c|}
\hline
时间段 & 购电量/kWh & 时间段 & 购电量/kWh & 时间段 & 购电量/kWh \\
\hline
10:00---10:10 & 0.00 & 12:00---12:10 & 480.41 & 14:00---14:10 & 0.00 \\
\hline
16:00---16:10 & 445.43 & 18:00---18:10 & 531.89 & 20:00---20:10 & 0.00 \\
\hline
\multicolumn{2}{|c|}{全天计划购电量/kWh} & 59,482.70 & \multicolumn{2}{c|}{全天购电费/元} & 35,126.95 \\
\hline
\end{tabular}
\end{table}

注：结果文件第一个数据位置对应 0:00---0:10，第 144 个数据位置对应
23:50---24:00；本表按实际时段取数。

\begin{table}[!htbp]
\centering
\caption{问题一指定时段充放电量及日初、日末储电量}
\small
\setlength{\tabcolsep}{4pt}
\begin{tabular}{|c|c|c|c|c|c|}
\hline
时间段 & 充电量/kWh & 放电量/kWh & 时间段 & 充电量/kWh & 放电量/kWh \\
\hline
0:00---4:00 & 4,500.00 & 0.00 & 4:00---8:00 & 833.33 & 6,365.84 \\
\hline
8:00---12:00 & 4,787.96 & 1,703.00 & 12:00---16:00 & 5,286.04 & 91.10 \\
\hline
16:00---20:00 & 0.00 & 5,780.13 & 20:00---24:00 & 5,333.33 & 2,859.87 \\
\hline
\multicolumn{2}{|c|}{0:00 储电量} & 6,000.00 & \multicolumn{2}{c|}{24:00 储电量} & 6,000.00 \\
\hline
\end{tabular}
\end{table}

注：四小时时段内充电量与放电量均为正，表示该时段内不同十分钟区间分别发生充电与放电，不表示同一时刻同时充放电。

\begin{figure}[!htbp]
\centering
\includegraphics[width=0.9\textwidth]{q1_01_plan_overview.png}
\caption{问题一供需平衡与储能充放电调度}
\end{figure}

由图 2
可见，外网购电、光伏出力与储能放电共同承担负载需求，储能充电使部分购电需求由高价时段转移至低价时段。全天累计充电量为
20,740.67 kWh，累计放电量为 16,799.94 kWh。两者满足

\begin{equation}Q_{D} = \eta_{c}\eta_{d}Q_{C} = 0.81Q_{C},\end{equation}

与日初、日末储电量相等时的能量关系一致，充放电量之差体现了储能转换损耗。

\begin{figure}[!htbp]
\centering
\includegraphics[width=0.85\textwidth]{q1_02_soc_price.png}
\caption{问题一储电量轨迹与分时电价}
\end{figure}

图 3
表明，储电量随充放电决策在运行区间内变化：低价阶段储能积累电量，高价阶段通过放电减少外购需求，并在日末恢复至初始状态。调度并非仅由电价决定，还受到光伏出力、负载需求及剩余储能容量的共同约束。

无储能基准购电费用为 48,052.05 元。与该基准相比，储能调度节省费用
12,925.10 元，降幅为 26.90\%。充电与放电的加权平均电价分别为 0.4930
元/kWh 和 1.1376 元/kWh，其比值约为
2.31，说明模型能够利用时段间价差，在承担能量转换损耗后降低购电费用。

\subsubsection{储能状态价值与模型验证}

为进一步解释储电量对购电成本的影响，在日末储电量固定为 6000 kWh
的条件下，以初始储电量 \(s\)为自变量，记最小购电费用为 \(V_{0}(s)\)。

\begin{figure}[!htbp]
\centering
\includegraphics[width=0.9\textwidth]{q1_08_value_function.png}
\caption{初始储电量对应的费用函数与边际价值}
\end{figure}

图4给出初始储电量与最小费用的关系。在 \(s = 6000\)kWh 附近，初始状态约束的对偶结果给出

\begin{equation}\frac{\partial J_{1}^{*}}{\partial s_{0}} = - 0.4754\text{~元}/kWh.\end{equation}

即在该局部范围内，初始储电量每增加 1 kWh，当天最优购电费用约减少 0.4754
元。

相应地，日末储电量约束的对偶结果为

\begin{equation}\frac{\partial J_{1}^{*}}{\partial s_{T}} = 0.4711\text{~元}/kWh,\end{equation}

表示日末留存要求每提高 1 kWh，当天最优购电费用约增加 0.4711
元。上述两个数值分别对应初始库存的利用价值与日末留存的边际成本，不能作为同一个量混用。

已有验证分别采用线性规划对偶、中心差分和离散动态规划检查初始库存的边际价值，结果相互一致。离散动态规划直接解的相对残差为
\(6.543 \times 10^{- 5}\)，经 Richardson 外推后为
\(2.524 \times 10^{- 6}\)；后者属于外推残差估计，不作为严格误差上界。结合能量平衡、储能状态递推及运行边界的独立复算，对调度结果进行交叉验证。

\subsection{问题二模型建立与求解}
\subsubsection{日前计划与实际执行的信息划分}

问题二采用``日前计划---实际执行''的两阶段框架。每天
0:00，在当日负载已知的假设下，利用此前已实现的光伏数据构造预测，制定全天购电计划；实际运行阶段，再根据真实光伏出力和储能状态执行调度。当天尚未实现的光伏数据不参与日前计划制定。

正式评价期为 2025 年 2 月 1 日至 12 月 31 日，共 334 天。1
月用于历史信息准备，储能采用中性运行方式，使 2 月 1 日 0:00
的初始储电量为 6000 kWh。此后，各日实际期末储电量连续传递至次日。

\begin{figure}[!htbp]
\centering
\includegraphics[width=0.95\textwidth]{q2_10_model_flow.png}
\caption{问题二日前计划与实际执行流程}
\end{figure}

如图 5
所示，历史信息用于生成计划，实际信息用于执行与结算，两类信息承担不同作用。计划购电费用按计划量计算，即使部分电量最终未被调用，也不能将普通购电费用改为按实际调用量结算。

\subsubsection{光伏预测与因果保守裕度}

设 \(E_{d,t}^{PV}\)为第 \(d\)日第
\(t\)时段的实际光伏电量，以过去四个连续自然日对应时段的均值作为基础预测：

\begin{equation}{\bar{E}}_{d,t}^{PV} = \frac{1}{4}\sum_{i = 1}^{4}E_{d - i,t}^{PV}.\end{equation}

对于基础预测为正的历史样本，定义相对预测误差

\begin{equation}r_{d,t} = \frac{E_{d,t}^{PV} - {\bar{E}}_{d,t}^{PV}}{{\bar{E}}_{d,t}^{PV}}.\end{equation}

利用决策日之前 28 天的历史相对误差确定保守裕度：

\begin{equation}\kappa_{d} = \mathrm{clip}\left(-Q_{0.2}\left(\mathcal{R}_{d}\right),\ 0,\ 0.35\right),\end{equation}

其中，\(\mathcal{R}_{d}\) 表示相应历史误差集合，\(Q_{0.2}\) 为其 20\%
分位数，\(clip\) 表示将数值限制在给定区间内。修正后的光伏预测为

\begin{equation}{\widehat{E}}_{d,t}^{PV} = \left( 1 - \kappa_{d} \right){\bar{E}}_{d,t}^{PV}.\end{equation}

该修正通过适当降低计划阶段采用的光伏电量，缓解光伏高估造成的供电缺口。分位水平用于体现普通购电与紧急购电的成本不对称，属于现有模型的启发式标定，不等同于包含储能耦合约束的全模型最优分位结论。裕度逐日利用历史信息更新，保持计划的非预见性。

\subsubsection{日前购电模型}

为区分计划与执行，分别使用上标 \(pl\)和 \(ex\)。日前计划变量为购电量
\(G_{d,t}^{pl}\)、充电量 \(C_{d,t}^{pl}\)、放电量
\(D_{d,t}^{pl}\)、弃光量 \(R_{d,t}^{pl}\)和储电量 \(S_{d,t}^{pl}\)。

计划层能量平衡为

\begin{equation}G_{d,t}^{pl} + {\widehat{E}}_{d,t}^{PV} + D_{d,t}^{pl} = E_{d,t}^{L} + C_{d,t}^{pl} + R_{d,t}^{pl}.\end{equation}

储能递推为

\begin{equation}S_{d,t}^{pl} = S_{d,t - 1}^{pl} + 0.9C_{d,t}^{pl} - \frac{D_{d,t}^{pl}}{0.9}.\end{equation}

购电、弃光与充放电变量满足

\begin{equation}{G_{d,t}^{pl} \geq 0,\ 0 \leq R_{d,t}^{pl} \leq {\widehat{E}}_{d,t}^{PV},
}{0 \leq C_{d,t}^{pl},D_{d,t}^{pl} \leq \frac{5000}{6},\ 1200 \leq S_{d,t}^{pl} \leq 10800.}\end{equation}

每天的计划初始状态取实际储电量：

\begin{equation}S_{d,0}^{pl} = S_{d,0}^{ex}.\end{equation}

为反映日末剩余电量的后续利用价值，采用已有的仿射终端价值，日前目标为

\begin{equation}\min J_{d}^{pl} = \sum_{t = 1}^{144}p_{t}G_{d,t}^{pl} - \lambda S_{d,144}^{pl},\ \lambda = 0.4720\text{~元}/kWh.\end{equation}

系数 \(\lambda=0.4720\) 元/kWh 依据问题一日末留存边际成本 \(0.4711\) 元/kWh 取近似标定值。该项只影响计划决策，不作为实际收入从现金费用中扣除；问题二也不强制各日日末储电量恢复至 6000 kWh。

\subsubsection{实际执行与费用核算}

计划购电按申报量 \(G_{d,t}^{pl}\) 结算，实际阶段再决定其中的调用量
\(U_{d,t}\)，并记未调用的已付费电量为
\[
W_{d,t}^{G}=G_{d,t}^{pl}-U_{d,t},\qquad 0\le U_{d,t}\le G_{d,t}^{pl}.
\]
令计划购电全部可用时的净缺口为
\[
g_{d,t}=E_{d,t}^{L}-E_{d,t}^{PV}-G_{d,t}^{pl}.
\]
若 \(g_{d,t}>0\)，先在功率和库存允许范围内放电，
\[
D_{d,t}^{ex}=\min\left\{g_{d,t},\frac{5000}{6},
0.9\left(S_{d,t-1}^{ex}-1200\right)\right\},\qquad
H_{d,t}=g_{d,t}-D_{d,t}^{ex},
\]
此时 \(U_{d,t}=G_{d,t}^{pl}\)。若 \(g_{d,t}\le0\)，则优先吸收富余电量，
\[
C_{d,t}^{ex}=\min\left\{-g_{d,t},\frac{5000}{6},
\frac{10800-S_{d,t-1}^{ex}}{0.9}\right\};
\]
超过负载与可充电量的部分依次表现为未调用计划电量或弃光。该执行顺序保证紧急购电仅用于储能放电后仍存在的真实缺口。各时段还满足
\[
U_{d,t}+H_{d,t}+E_{d,t}^{PV}+D_{d,t}^{ex}
=E_{d,t}^{L}+C_{d,t}^{ex}+R_{d,t}^{PV},
\]
\[
S_{d,t}^{ex}=S_{d,t-1}^{ex}+0.9C_{d,t}^{ex}
-\frac{D_{d,t}^{ex}}{0.9},\qquad
S_{d+1,0}^{ex}=S_{d,144}^{ex},
\]
\[
1200\le S_{d,t}^{ex}\le10800,\quad
0\le C_{d,t}^{ex},D_{d,t}^{ex}\le\frac{5000}{6},\quad
0\le R_{d,t}^{PV}\le E_{d,t}^{PV}.
\]
其中 \(W_{d,t}^{G}\) 与弃光量 \(R_{d,t}^{PV}\) 含义不同：前者已经付费但未调用，后者是未利用的光伏电量。

评价期现金费用仅由计划购电费与紧急购电费构成：
\[
J_{2}=\underbrace{\sum_{d,t}p_tG_{d,t}^{pl}}_{J_{plan}}
+\underbrace{\sum_{d,t}5p_tH_{d,t}}_{J_{emg}}.
\]
终端价值只影响日前决策，不从现金费用中扣除。由于紧急购电单价为普通电价的五倍，降低光伏预测会提高普通购电计划，以较小的冗余成本换取对高价缺口的抑制，而不是减少计划购电。
\subsubsection{指定日期结果}

为减少重复表头并控制篇幅，四个指定日期采用同口径合并列示。表中购电量为日前计划量
\(G_{d,t}^{pl}\)，储能量为执行层结果；各时刻列均对应其后十分钟区间。

\begin{table}[!htbp]
\centering\scriptsize
\caption{问题二指定日期的计划购电量（kWh）}
\setlength{\tabcolsep}{3pt}
\resizebox{\textwidth}{!}{
\begin{tabular}{|c|r|r|r|r|r|r|r|}
\hline
日期 & 10:00 & 12:00 & 14:00 & 16:00 & 18:00 & 20:00 & 全天合计 \\
\hline
2025年3月20日 & 0.00 & 552.31 & 0.00 & 519.50 & 698.00 & 0.00 & 67,083.45 \\
\hline
2025年6月21日 & 0.00 & 0.00 & 0.00 & 185.94 & 0.00 & 0.00 & 34,795.81 \\
\hline
2025年9月23日 & 0.00 & 555.11 & 0.00 & 623.07 & 764.33 & 0.00 & 70,722.56 \\
\hline
2025年12月21日 & 0.00 & 1,002.77 & 0.00 & 806.43 & 715.41 & 0.00 & 95,876.74 \\
\hline
\end{tabular}}
\end{table}

\begin{table}[!htbp]
\centering\small
\caption{问题二指定日期的购电费用（元）}
\setlength{\tabcolsep}{6pt}
\begin{tabular}{|c|r|r|r|}
\hline
日期 & 计划购电费 & 紧急购电费 & 当日合计 \\
\hline
2025年3月20日 & 40,498.51 & 1,096.59 & 41,595.10 \\
\hline
2025年6月21日 & 19,334.88 & 0.00 & 19,334.88 \\
\hline
2025年9月23日 & 43,914.08 & 181.89 & 44,095.97 \\
\hline
2025年12月21日 & 61,446.07 & 0.00 & 61,446.07 \\
\hline
\end{tabular}
\end{table}

\begin{table}[!htbp]
\centering\scriptsize
\caption{问题二指定日期的储能充放电量与日界储电量（kWh）}
\setlength{\tabcolsep}{2pt}
\resizebox{\textwidth}{!}{
\begin{tabular}{|c|c|c|c|c|c|c|r|r|}
\hline
日期 & 0:00---4:00 & 4:00---8:00 & 8:00---12:00 & 12:00---16:00 & 16:00---20:00 & 20:00---24:00 & \(S_{0}\) & \(S_{24}\) \\
\hline
 & \multicolumn{6}{c|}{各区间充电量/放电量} & & \\
\hline
2025年3月20日 & 2,332.84/0.00 & 1,010.17/5,526.08 & 5,937.65/2,777.38 & 4,451.16/364.61 & 41.73/5,790.89 & 7,500.00/2,773.02 & 7,950.44 & 7,950.00 \\
\hline
2025年6月21日 & 0.00/1,311.73 & 344.45/4,405.91 & 9,849.77/0.00 & 0.00/0.00 & 16.05/6,127.58 & 7,500.00/2,485.90 & 7,978.14 & 7,993.91 \\
\hline
2025年9月23日 & 2,333.33/0.00 & 889.56/5,635.92 & 6,293.08/1,896.62 & 3,525.40/465.97 & 0.93/5,629.73 & 7,500.00/3,011.02 & 7,950.00 & 7,950.00 \\
\hline
2025年12月21日 & 2,333.32/0.00 & 1,666.66/1,483.03 & 5,855.25/7,112.37 & 6,663.91/2,223.29 & 3.93/5,797.47 & 7,500.00/2,842.53 & 7,950.01 & 7,950.00 \\
\hline
\end{tabular}}
\end{table}

\begin{table}[!htbp]
\centering\small
\caption{问题二指定日期的紧急购电}
\begin{tabular}{|c|c|c|}
\hline
日期 & 时间段 & 紧急购电量/kWh \\
\hline
2025年3月20日 & 20:30---20:40 & 157.19 \\
\hline
2025年6月21日 & 无 & 0.00 \\
\hline
2025年9月23日 & 20:30---20:40 & 26.07 \\
\hline
2025年12月21日 & 无 & 0.00 \\
\hline
\end{tabular}
\end{table}

四个指定日期中，6月21日的全天计划购电量与费用最低，12月21日最高；紧急购电只在
3月20日和9月23日的20:30---20:40发生。四小时区间内充、放电量同时为正，只表示不同十分钟时段分别充电和放电，并非同一时刻同时进行。
\subsubsection{评价期费用、对照与验证}

评价期334天的计划购电费为 \(12{,}891{,}818.24\) 元，紧急购电费为
\(360{,}522.85\) 元，故
\[
J_2=12{,}891{,}818.24+360{,}522.85
=13{,}252{,}341.09\text{ 元}.
\]
两部分分别占总费用的97.28\%和2.72\%。

\begin{figure}[!htbp]
\centering
\includegraphics[width=0.9\textwidth]{q2_04_cost_structure.png}
\caption{问题二评价期购电费用结构与方案对照}
\end{figure}

图6给出评价期费用构成与方案对照。评价期累计计划购电量为 \(21{,}283{,}432.62\) kWh，紧急购电量为
\(55{,}801.05\) kWh；紧急购电发生于242天，共形成375个合并区间。累计充电量与放电量分别为
\(6{,}684{,}044.85\) kWh 和 \(5{,}412{,}321.33\) kWh，储电量由
6000 kWh 变为约7950 kWh，并满足
\[
S_{end}-S_{start}=0.9\sum_{d,t}C_{d,t}^{ex}
-\frac{1}{0.9}\sum_{d,t}D_{d,t}^{ex}.
\]
因此，储能损耗、期末库存变化与供电缺口可以在统一能量账本中区分。

\begin{table}[!htbp]
\centering\footnotesize
\caption{问题二预报与裕度方案的统一对照（全年现金费用）}
\setlength{\tabcolsep}{3pt}
\resizebox{\textwidth}{!}{%
\begin{tabular}{|l|r|r|l|}
\hline
方案 & 总费用/元 & 相对主方案 & 对照含义 \\
\hline
\textbf{本文主方案} & \textbf{13,252,341.09} & --- & 四日均值点预测与逐日因果裕度 \\
\hline
同一预测、裕度为0 & 14,178,004.00 & +925,662.91（+6.98\%） & 同口径隔离裕度作用 \\
\hline
四情景SAA、裕度为0 & 13,396,023.55 & +143,682.46（+1.08\%） & 不同预测方法，仅作方法参照 \\
\hline
完全信息下界 & 12,254,696.55 & $-$997,644.54（$-$7.53\%） & 当日实际光伏预先已知，不可在线实施 \\
\hline
同周四周情景 & 15,238,244.28 & +1,985,903.19（+14.99\%） & 信息受限对照 \\
\hline
附件1典型日曲线 & 19,043,406.75 & +5,791,065.66（+43.70\%） & 固定典型预测对照 \\
\hline
\end{tabular}}
\end{table}

与同一预测且不设裕度的方案相比，主方案节省
\(925{,}662.91\) 元，按无裕度方案计降幅为6.53\%，说明因果裕度能够以增加普通计划购电来抑制五倍价格的紧急购电。四情景SAA对照改变了预测方法，不能据其与主方案的费用差单独归因于裕度。完全信息方案仅给出理论下界；主方案比该下界高
\(997{,}644.54\) 元，占主方案费用的7.53\%，该差额反映信息受限与保守决策的综合代价。

终端价值系数在不高于0.40元/kWh时总费用约为
\(13{,}251{,}225.13\) 元，取0.4720元/kWh时为本文结果，取值不低于0.55元/kWh时约为
\(13{,}287{,}931.82\) 元；区间极差为 \(36{,}706.69\) 元，仅占主方案费用约0.28\%。这表明结果对该系数存在响应，但总体费用排序较稳定。

最后按“数据位置—十分钟区间”映射独立汇总各费用项，并逐时检查能量平衡、储能递推、功率和容量边界、跨日状态衔接以及
\(J_2=J_{plan}+J_{emg}\)。同时将决策日之后的光伏数据隔离于预测构造之外，以检验日前计划的非预见性。上述复核与主结果一致。
\subsection{问题三模型建立与求解}

\subsubsection{预报更新与决策过程}

问题三在日前计划的基础上增加光伏预报更新与购电调整机制。每天 0:00
制定初始计划，随后利用 6:00、12:00 和 18:00
发布的预报修正未来时段的购电安排，实际运行时结合供需情况与储电量执行调度。各阶段共享储能状态，但承担不同决策任务：计划层确定初始购电量，调整层修正尚未执行的策略，执行层处理预测与实际之间的偏差。

\begin{figure}[!htbp]
\centering
\includegraphics[width=0.9\textwidth]{q3_16_rollout_timeline.png}
\caption{光伏预报更新与购电调整时序}
\end{figure}

如图 7
所示，每次预报仅影响发布时刻之后的决策，已经执行的购电与充放电动作不再改变。同一时段可能经历多次策略更新，但费用仅根据其最终生效的调整量相对
0:00 计划量结算，避免对中间调整版本重复计费。

附件 3 提供整点光伏预报，需要转换为十分钟尺度。设发布时刻为
\(\tau\)，相邻预报节点为 \(F_{k},F_{k + 1}\)，对应时刻为
\(u_{k},u_{k + 1}\)，则节点间采用线性插值：

\begin{equation}{\widehat{P}}_{\tau}^{PV}(u) = F_{k} + \frac{u - u_{k}}{u_{k + 1} - u_{k}}\left( F_{k + 1} - F_{k} \right),\ u \in \left\lbrack u_{k},u_{k + 1} \right\rbrack.\end{equation}

首个区间的左侧锚点取发布时刻可获得的已实现光伏信息。以小时为积分单位，将重建功率转换为时段电量：

\begin{equation}{\widehat{E}}_{\tau,t}^{PV} = \int_{u_{t - 1}}^{u_{t}}{\widehat{P}}_{\tau}^{PV}(u)\text{ }du.\end{equation}

由此使光伏预报与购电、储能调度采用一致的时间尺度。

\subsubsection{计划量、调整量与费用结算}

设 \(G_{d,t}^{0}\)为第 \(d\)日 0:00 制定的计划购电量，\(A_{d,t}\) 为第
\(t\)时段最终生效的调整购电量，\(H_{d,t}\) 为紧急购电量。记

\begin{equation}(x)^{+} = \max\left( x,0 \right).\end{equation}

按照本文对题面偏差计价规则的解释，单时段现金费用为

\begin{equation}\begin{aligned}
c_{d,t} = & p_{t}\min\left( G_{d,t}^{0},A_{d,t} \right) \\
 & + 0.5p_{t}\left( G_{d,t}^{0} - A_{d,t} \right)^{+} \\
 & + 1.5p_{t}\left( A_{d,t} - G_{d,t}^{0} \right)^{+} + 5p_{t}H_{d,t}.
\end{aligned}\end{equation}

其中，保留的计划电量按正常电价结算，取消的计划电量按正常电价的 50\%
计费，新增电量按正常电价的 1.5 倍计费，供电缺口另按 5 倍电价紧急购电。

为与``计划购电费用、调整相关费用及紧急购电费用''对应，上式亦可写为

\begin{equation}c_{d,t} = p_{t}G_{d,t}^{0} - 0.5p_{t}\left( G_{d,t}^{0} - A_{d,t} \right)^{+} + 1.5p_{t}\left( A_{d,t} - G_{d,t}^{0} \right)^{+} + 5p_{t}H_{d,t}.\end{equation}

此时，调整相关费用是相对原计划费用的增减额，下调对应净退还一半费用。两种表达完全等价，但不能将第一种表达中的完整普通购电结算额再与原计划费用相加。

该结算口径下，单纯减少购电量并不必然降低总费用：若调整后供电不足，节省的普通购电费用可能被紧急购电费用抵消。因此，调整决策还需结合光伏预测与储能调节能力。

\subsubsection{滚动调度与储能约束}

计划阶段沿用已有的净需求保守度参数 \(q\)与光伏相关裕度参数
\(m\)的参数化构造。二者共同影响计划输入及购电决策，属于参数标定对象，不能仅凭不同配置的费用差判定某一分位方法具有普遍优势。

每次更新后，购电与储能安排均满足供需平衡和运行边界。储能状态方程保持为

\begin{equation}S_{d,t} = S_{d,t - 1} + 0.9C_{d,t} - \frac{D_{d,t}}{0.9},\end{equation}

并满足

\begin{equation}1200 \leq S_{d,t} \leq 10800,\ 0 \leq C_{d,t},D_{d,t} \leq \frac{5000}{6}.\end{equation}

实际执行中，设 \(U_{d,t}\)为实际调用的普通购电量，则

\begin{equation}{0 \leq U_{d,t} \leq A_{d,t},
}{U_{d,t} + H_{d,t} + E_{d,t}^{PV} + D_{d,t} = E_{d,t}^{L} + C_{d,t} + R_{d,t}.}\end{equation}

其中 \(R_{d,t}\)为弃光量，实际充放电满足互斥条件

\begin{equation}C_{d,t}D_{d,t} = 0.\end{equation}

不同决策时刻通过实际储电量衔接，避免后续优化使用已经偏离实际运行的计划状态。

日末剩余电量采用既有终端价值项

\begin{equation}- \lambda_{T}S_{d,144},\ \lambda_{T} = 0.4720\text{~元}/kWh\end{equation}

计入优化目标，以体现储电量的后续利用价值。该项用于指导调度，评价期现金费用仍按前述结算公式累计，不将终端价值作为实际收入扣除。

实际供需与原定充放电动作冲突时，执行层允许撤销无法由实际供需支持的计划动作，再依据当前储电量重新分配充放电；仅在可用储能放电后仍有真实缺口时紧急购电。由此保证执行序列物理可行，并避免为完成不再合理的计划充电而额外购买高价紧急电。

\subsubsection{指定日期购电与储能结果}

四个指定日期统一采用相同口径：计划量 \(G_{d,t}^{0}\) 为0:00确定的原始购电计划，调整量
\(A_{d,t}\) 为各次预报更新后最终生效的购电量；各时刻列均对应其后十分钟区间。将同构结果合并列示，以减少重复表头。

\begin{table}[!htbp]
\centering\scriptsize
\caption{问题三指定日期的计划与调整购电量（kWh）}
\setlength{\tabcolsep}{2.5pt}
\resizebox{\textwidth}{!}{%
\begin{tabular}{|c|c|r|r|r|r|r|r|r|}
\hline
日期 & 类型 & 10:00 & 12:00 & 14:00 & 16:00 & 18:00 & 20:00 & 全天合计 \\
\hline
\multirow{2}{*}{2025年3月20日} & 计划 & 0.00 & 621.87 & 0.00 & 692.13 & 698.00 & 0.00 & 70,918.39 \\
\cline{2-9}
 & 调整 & 0.00 & 491.43 & 0.00 & 543.53 & 698.00 & 0.00 & 65,979.69 \\
\hline
\multirow{2}{*}{2025年6月21日} & 计划 & 0.00 & 0.00 & 0.00 & 156.76 & 0.00 & 0.00 & 34,499.06 \\
\cline{2-9}
 & 调整 & 0.00 & 0.00 & 0.00 & 102.64 & 0.00 & 0.00 & 34,173.10 \\
\hline
\multirow{2}{*}{2025年9月23日} & 计划 & 0.00 & 452.48 & 0.00 & 694.96 & 792.73 & 0.00 & 73,703.11 \\
\cline{2-9}
 & 调整 & 0.00 & 0.00 & 0.00 & 597.11 & 793.54 & 0.00 & 67,644.50 \\
\hline
\multirow{2}{*}{2025年12月21日} & 计划 & 0.00 & 1,069.53 & 0.00 & 839.91 & 715.41 & 0.00 & 96,785.58 \\
\cline{2-9}
 & 调整 & 0.00 & 878.64 & 0.00 & 839.91 & 715.41 & 0.00 & 94,515.18 \\
\hline
\end{tabular}}
\end{table}

\begin{table}[!htbp]
\centering\scriptsize
\caption{问题三指定日期的现金费用分解（元）}
\setlength{\tabcolsep}{4pt}
\begin{tabular}{|c|r|r|r|r|r|}
\hline
日期 & 保留购电结算 & 下调偏差费 & 上调购电费 & 紧急购电费 & 当日合计 \\
\hline
2025年3月20日 & 39,804.28 & 1,692.14 & 374.95 & 1,839.82 & 43,711.19 \\
\hline
2025年6月21日 & 18,784.09 & 142.43 & 15.61 & 13.05 & 18,955.18 \\
\hline
2025年9月23日 & 42,388.71 & 1,887.59 & 17.08 & 244.25 & 44,537.63 \\
\hline
2025年12月21日 & 60,028.12 & 961.02 & 625.42 & 560.57 & 62,175.13 \\
\hline
\end{tabular}
\end{table}

表中四项费用依次对应
\(p_t\min(G_{d,t}^{0},A_{d,t})\)、
\(0.5p_t(G_{d,t}^{0}-A_{d,t})^{+}\)、
\(1.5p_t(A_{d,t}-G_{d,t}^{0})^{+}\) 和 \(5p_tH_{d,t}\)，其和为当日现金费用。

\begin{table}[!htbp]
\centering\scriptsize
\caption{问题三指定日期的储能充放电量与日界储电量（kWh）}
\setlength{\tabcolsep}{2pt}
\resizebox{\textwidth}{!}{%
\begin{tabular}{|c|c|c|c|c|c|c|r|r|}
\hline
日期 & 0:00---4:00 & 4:00---8:00 & 8:00---12:00 & 12:00---16:00 & 16:00---20:00 & 20:00---24:00 & \(S_0\) & \(S_{24}\) \\
\hline
 & \multicolumn{6}{c|}{各区间充电量/放电量} & & \\
\hline
2025年3月20日 & 2,332.27/0.00 & 841.07/6,073.62 & 4,478.39/2,468.26 & 6,083.54/248.21 & 197.69/5,640.41 & 7,500.00/2,930.22 & 7,950.95 & 7,950.92 \\
\hline
2025年6月21日 & \(<0.01\)/1,333.46 & 399.38/4,389.23 & 9,809.90/0.00 & 0.00/0.00 & 24.66/6,090.47 & 7,500.00/2,485.90 & 7,970.19 & 8,042.90 \\
\hline
2025年9月23日 & 2,333.33/0.00 & 843.92/6,343.53 & 6,128.07/1,895.25 & 4,491.03/371.26 & 0.00/5,602.90 & 7,500.00/3,037.09 & 7,950.01 & 7,950.00 \\
\hline
2025年12月21日 & 2,330.27/0.00 & 1,664.10/1,506.25 & 5,269.27/7,706.18 & 8,014.21/2,220.11 & 0.00/5,797.47 & 7,500.00/2,842.53 & 7,952.76 & 7,950.00 \\
\hline
\end{tabular}}
\end{table}

四小时区间内充、放电量同时为正只表示不同十分钟时段分别充电和放电；“\( <0.01 \)”表示非零但不足0.01 kWh。

\begin{table}[!htbp]
\centering\scriptsize
\caption{问题三指定日期的紧急购电}
\setlength{\tabcolsep}{2.5pt}
\resizebox{\textwidth}{!}{%
\begin{tabular}{|c|r|c|r|c|r|c|r|}
\hline
\multicolumn{2}{|c|}{2025年3月20日} & \multicolumn{2}{c|}{2025年6月21日} & \multicolumn{2}{c|}{2025年9月23日} & \multicolumn{2}{c|}{2025年12月21日} \\
\hline
时间段 & kWh & 时间段 & kWh & 时间段 & kWh & 时间段 & kWh \\
\hline
6:40---7:00 & 4.88 & 4:00---4:50 & 1.40 & 6:00---6:20 & 9.89 & 7:00---7:30 & 3.99 \\
\hline
7:40---8:40 & 91.62 & 19:00---19:50 & 1.48 & 10:00---10:10 & 1.36 & 7:50---8:00 & 2.11 \\
\hline
8:50---9:20 & 151.22 & --- & 0.00 & 14:20---14:50 & 32.36 & 9:00---9:20 & 4.67 \\
\hline
9:30---10:00 & 77.61 & --- & 0.00 & 17:30---17:50 & 5.45 & 9:30---10:50 & 108.27 \\
\hline
14:10---14:20 & 6.51 & --- & 0.00 & --- & 0.00 & --- & 0.00 \\
\hline
15:30---15:40 & 0.72 & --- & 0.00 & --- & 0.00 & --- & 0.00 \\
\hline
15:50---16:20 & 33.78 & --- & 0.00 & --- & 0.00 & --- & 0.00 \\
\hline
\end{tabular}}
\end{table}

滚动更新后仍可能产生紧急购电，且区间不局限于晚高峰；是否出现缺口同时取决于预报误差、储能状态和最终生效购电量。
\subsubsection{评价期费用与储能运行}

在 \(q=0.8\)、\(m=0\) 和334天评价期口径下，采用上述允许回退冲突动作的执行规则，评价期总现金费用为

\begin{equation}J_{3} = 13{,}120{,}194.06\text{~元},\end{equation}

累计紧急购电量为

\begin{equation}\sum_{d,t}^{}H_{d,t} = 42{,}672.22\text{~}kWh.\end{equation}

该结果对应全文统一的信息权限、执行规则与结算口径。

按照既有结果中的结算分块，保留计划电量对应费用为 12,381,239.55
元，取消及新增电量对应费用合计为 542,052.30 元，紧急购电费用为 196,902.21
元，三者相加得到上述总费用。这里第一项对应

\begin{equation}\sum_{d,t}^{}p_{t}\min\left( G_{d,t}^{0},A_{d,t} \right),\end{equation}

不等同于原始计划量的名义购电费用。

上述费用分解将普通购电结算与紧急购电区分开来，使策略评价同时考虑调整成本和实际供电缺口。单独压低某一费用项并不代表总费用更低，应以完整现金费用作为经济性评价依据。

\begin{figure}[!htbp]
\centering
\includegraphics[width=0.9\textwidth]{q3_05_soc_year.png}
\caption{问题三评价期储电量轨迹与运行边界}
\end{figure}

储能在各日间连续运行，评价期初始储电量为 6000 kWh，期末储电量约为
7950.00 kWh。图 8
同时展示日内储电量范围与跨日状态，说明仅考察日初或日末储电量不足以判断整个运行过程是否满足容量约束，需结合逐时段状态进行验证。

\subsubsection{不同预报时刻的信息价值}

为分析是否需要引入其他时刻的预报，采用已有的四组预报更新对照：

\begin{equation}\begin{aligned}
\mathcal{A}_{0} & = \left\{ 0:00 \right\}, \\
\mathcal{A}_{1} & = \left\{ 0{:}00,\ 6{:}00 \right\}, \\
\mathcal{A}_{2} & = \left\{ 0{:}00,\ 6{:}00,\ 12{:}00 \right\}, \\
\mathcal{A}_{3} & = \left\{ 0{:}00,\ 6{:}00,\ 12{:}00,\ 18{:}00 \right\}.
\end{aligned}\end{equation}

相应评价期费用如下。

\begin{table}[!htbp]
\centering
\caption{问题三滚动信息价值（四组预报更新方案）}
\begin{tabular}{|c|c|c|}
\hline
预报更新方案 & 全年费用/元 & 相对上一方案节省/元 \\
\hline
仅 0:00 & 13,564,151.84 & --- \\
\hline
增加 6:00 & 13,446,385.03 & 117,766.81 \\
\hline
再增加 12:00 & 13,125,320.04 & 321,064.99 \\
\hline
再增加 18:00 & 13,120,194.06 & 5,125.98 \\
\hline
\end{tabular}
\end{table}

\begin{figure}[!htbp]
\centering
\includegraphics[width=0.9\textwidth]{q3_03_rolling_value_voi.png}
\caption{不同预报更新时刻的费用变化与边际贡献}
\end{figure}

由图9与上表可见，在逐步加入预报的比较顺序下，6:00 与 12:00
更新均带来费用下降，其中加入 12:00 预报的增量收益最大；进一步加入 18:00
预报仍可降低费用，但收益明显较小。完整更新组合相对仅使用 0:00
预报，费用减少 443,957.78 元，降幅约为 3.27\%。

这说明，在现有数据与策略条件下，引入日内预报调整具有经济价值，尤其应重视白天新增信息对后续供需判断的修正。18:00
更新的增量收益较小，但不为零。上述结论适用于已给定的更新时刻及其组合，不能据此直接推断任意增加预报频率都能获得相同收益。

\subsection{问题四模型建立与求解}

\subsubsection{波动电价下的模型扩展}

问题四将固定日电价 \(p_{t}\)扩展为随日期和时段变化的电价
\(p_{d,t}\)。依据本文假设，每天 0:00 可获得附件 4
给定的当日电价序列，并据此制定购电策略。光伏信息、储能参数及相应交易规则仍分别沿用问题二和问题三。

\begin{figure}[!htbp]
\centering
\includegraphics[width=0.9\textwidth]{q4_01_price_caliber.png}
\caption{典型日电价与波动电价的时序特征}
\end{figure}

图 10
展示了同一时段在不同日期的价格分布及日均价格变化。固定日电价能够描述典型日内变化，但不能反映各日价格水平的差异。价格变化会改变购电与储能充放电的经济排序，因此，需要在新价格条件下重新求解完整策略，而非仅对原有购电量更换价格后计算费用。

本问保留前述能量平衡、储能状态递推、容量边界及充放电限制。两条计算分支分别对应波动电价下的日前计划策略与滚动调整策略。

\subsubsection{波动电价下的日前计划策略}

问题四对应问题二的分支继续采用历史光伏均值与因果保守裕度生成预测，按每日已知价格确定购电和储能计划。日前目标为

\begin{equation}\min J_{d,4\text{-}2}^{pl} = \sum_{t = 1}^{144}p_{d,t}G_{d,t}^{pl} - \lambda_{T}S_{d,144}^{pl},\ \lambda_{T} = 0.4720\text{~元}/kWh.\end{equation}

实际运行采用问题二的执行规则，普通购电按计划量结算，供电缺口按当前电价的
5 倍紧急购电。评价期现金费用为

\begin{equation}J_{4\text{-}2} = \sum_{d}^{}{\sum_{t}^{}p_{d,t}}G_{d,t}^{pl} + \sum_{d}^{}{\sum_{t}^{}5}p_{d,t}H_{d,t}.\end{equation}

终端价值项仅用于优化决策，不计入上述现金费用。

该分支保持原有预测方法及执行机制，使结果能够反映价格条件变化对同一类策略的影响。各日实际期末储电量继续传递至次日，不重新设置每日初始状态。

\subsubsection{波动电价下的滚动调整策略}

问题四对应问题三的分支完整继承0:00原始计划、6:00/12:00/18:00预报更新、净结算口径及允许回退冲突动作的执行规则。各次更新只调整尚未执行的时段，费用按最终生效调整量相对原计划结算。

将问题三的固定价格替换为 \(p_{d,t}\)，总现金费用为

\begin{equation}\begin{aligned}
J_{4\text{-}3} = \sum_{d}^{}{\sum_{t}^{}\lbrack} & p_{d,t}\min\left( G_{d,t}^{0},A_{d,t} \right) \\
 & + 0.5p_{d,t}\left( G_{d,t}^{0} - A_{d,t} \right)^{+} \\
 & + 1.5p_{d,t}\left( A_{d,t} - G_{d,t}^{0} \right)^{+} + 5p_{d,t}H_{d,t}\rbrack.
\end{aligned}\end{equation}

该式已包含完整普通购电结算与紧急购电费用，无须再次叠加原计划购电费用。

与问题三相同，实际调用的普通购电量不超过最终生效调整量，充放电满足运行边界和互斥要求。价格变化同时影响计划购电、调整成本及紧急购电成本，因此，两条分支均在波动电价下重新执行完整计算。

\subsubsection{评价期计算结果}

2025 年 2 月 1 日至 12 月 31
日的计算结果如下。固定电价结果分别取问题二和问题三的对应方案，其中Q4-3的基线为问题三滚动调整结果。

\begin{table}[!htbp]
\centering
\caption{问题四波动电价下的全年结果与归因}
\small
\setlength{\tabcolsep}{4pt}
\begin{tabular}{|l|r|r|}
\hline
指标 & Q4-2（日前计划策略） & Q4-3（滚动调整策略） \\
\hline
全年费用/元 & 13,850,454.64 & 13,727,033.66 \\
\hline
本问基线/元 & 13,252,341.09（问题二） & 13,120,194.06（问题三） \\
\hline
相对基线/元 & 598,113.55 & 606,839.60 \\
\hline
价格水平效应/元 & 744,522.48 & 760,032.97 \\
\hline
日前重优化效应/元 & -146,408.94 & -153,193.38 \\
\hline
两项合计/元 & 598,113.55 & 606,839.60 \\
\hline
价格冲击占比 & 5.62\% & 5.79\% \\
\hline
重优化对冲比例 & 19.66\% & 20.16\% \\
\hline
\end{tabular}
\vspace{2pt}
\begin{flushleft}\footnotesize 注：价格水平效应为“同策略换价”，日前重优化效应为“同价换策略”；两组基线分属不同问题，组内可比，组间仅比较效应量级。分项按未舍入值相加，显示值可能存在 0.01 元舍入差。\end{flushleft}
\par\smallskip
\textbf{Q4-3滚动分支交叉结算与机制分解}\par\smallskip
\begin{tabular}{|c|r|r|}
\hline
计划来源/结算价格 & 附件1固定价 & 附件4波动价 \\
\hline
附件1价格下优化 & 13,120,194.06 & 13,880,227.04 \\
\hline
附件4价格下优化 & 13,225,150.38 & 13,727,033.66 \\
\hline
效应 & 数值/元 & 含义 \\
\hline
价格环境效应（旧计划） & +760,032.97 & 同一计划换价结算 \\
\hline
重新优化效应（新价格） & $-$153,193.38 & 同一价格更换计划 \\
\hline
合计 & +606,839.60 & Q4-3相对Q3（+4.63\%） \\
\hline
反向：旧价格下重新优化 & +104,956.32 & 机制诊断 \\
\hline
反向：新计划的价格环境 & +501,883.28 & 机制诊断 \\
\hline
差上差交互 & $-$258,149.70 & 不解释为算法收益 \\
\hline
分解残差 & 0.00 & 按未舍入值计算 \\
\hline
\end{tabular}\end{table}

\begin{table}[!htbp]
\centering
\caption{问题四指定日期的全天购电量与购电费}
\footnotesize
\setlength{\tabcolsep}{4pt}
\begin{tabular}{|l|r|r|r|r|}
\hline
日期（2025 年） & Q4-2 购电量/kWh & Q4-2 购电费/元 & Q4-3 购电量/kWh & Q4-3 购电费/元 \\
\hline
3 月 20 日 & 67,948.54 & 42,519.33 & 71,794.22 & 44,691.48 \\
\hline
6 月 21 日 & 35,601.06 & 16,794.81 & 35,317.51 & 16,549.17 \\
\hline
9 月 23 日 & 69,050.22 & 44,724.64 & 71,846.20 & 45,275.49 \\
\hline
12 月 21 日 & 90,210.08 & 69,595.07 & 91,121.97 & 70,593.86 \\
\hline
四日期合计 & 262,809.90 & 173,633.85 & 270,079.90 & 177,110.00 \\
\hline
全年合计（334 天） & 21,323,915.22 & 13,850,454.64 & 21,541,336.55 & 13,727,033.66 \\
\hline
\end{tabular}
\vspace{2pt}
\begin{flushleft}\footnotesize 注：末行为 334 天评价期合计；上方“四日期合计”仅包含题目点名的四个日期。Q4-2 购电费为计划购电费与紧急购电费之和，Q4-3为本文净结算口径下的当日费用。\end{flushleft}
\end{table}


四个指定日期及 334 天评价期的购电量、购电费均列于上表。两类策略均通过储能和购电计划调整参与供需协调；其中 Q4-2 购电费为计划购电费与紧急购电费之和，Q4-3购电费为本文净结算口径下的当日费用。

\begin{figure}[!htbp]
\centering
\includegraphics[width=0.95\textwidth]{q4_02_plan_heatmap.png}
\caption{波动电价下 Q4-2 计划购电量的日内分布与逐日合计}
\end{figure}

图11显示，计划购电量在334天内呈现明显的日内与日际差异。热力图反映了负载、光伏及价格共同作用下的购电时序结构，逐日合计则展示不同日期总外购需求的变化。

\begin{figure}[!htbp]
\centering
\includegraphics[width=0.95\textwidth]{q4_03_soc_price_7d.png}
\caption{波动电价下代表日期的储电量与电价变化}
\end{figure}

图 12
展示了代表日期的储电量轨迹。储能在部分低价时段积累电量，并在后续需求较高的时段放电，但具体动作还受光伏富余与剩余容量限制。因此，价格变化主要通过改变电量在各时段的配置影响费用，不能仅用``低价充电、高价放电''概括全部运行行为。

\subsubsection{价格环境与策略调整的费用分解}

仅比较重新优化前后的总费用，不能区分成本变化来自价格环境还是策略响应。为此，采用已有的
\(2 \times 2\)交叉结算结果，分别记录两种价格条件下生成的策略在两种结算价格下的费用。

记固定电价为 \(p^{(0)}\)，波动电价为 \(p^{(1)}\)，相应生成的完整策略为
\(\Pi^{(0)}\)和 \(\Pi^{(1)}\)。以

\begin{equation}\mathcal{C}\left( \Pi,p \right)\end{equation}

表示给定策略在价格
\(p\)下的现金费用。沿``先更换结算价格、再重新优化策略''的分解顺序，有

\begin{equation}\begin{aligned}
\Delta J = & \mathcal{C}\left( \Pi^{(1)},p^{(1)} \right)-\mathcal{C}\left( \Pi^{(0)},p^{(0)} \right) \\
 = {} & \underbrace{\mathcal{C}\left( \Pi^{(0)},p^{(1)} \right)-\mathcal{C}\left( \Pi^{(0)},p^{(0)} \right)}_{\Delta J_{price}} \\
 & + \underbrace{\mathcal{C}\left( \Pi^{(1)},p^{(1)} \right)-\mathcal{C}\left( \Pi^{(0)},p^{(1)} \right)}_{\Delta J_{strategy}}.
\end{aligned}\end{equation}

其中，第一项保持原策略不变，衡量价格环境变化的费用影响；第二项保持波动结算价格不变，衡量重新优化策略后的费用变化。


\begin{figure}[!htbp]
\centering
\includegraphics[width=0.9\textwidth]{q4_05_cf_2x2.png}
\caption{价格环境变化与策略重优化的交叉结算对照}
\end{figure}

图13中，Q4-2 与 Q4-3 分别以问题二和问题三滚动调整结果为各自基线，
因此只在各自组内比较价格效应与重优化效应，不将两组格值直接横向作差。

按照已有分解结果，日前计划策略的价格环境效应约为增加 74.45
万元，重新优化策略抵消约 14.64 万元；滚动调整策略的价格环境效应约为增加
76.00 万元，重新优化策略抵消约 15.32
万元。两类策略均通过重新配置购电与储能运行缓解了价格冲击，但节省额不足以完全抵消外部价格变化带来的成本增加。

\begin{figure}[!htbp]
\centering
\includegraphics[width=0.85\textwidth]{q4_06_effect_split.png}
\caption{日前计划策略的分时段价格影响与重优化效果}
\end{figure}

图 14
进一步展示了日前计划策略的时段差异。重新优化后，部分夜间时段支出增加，而午间及晚峰时段支出减少，最终形成净节省。因而策略响应并非要求每个时段均减少购电，而是允许部分时段增加支出，以换取其他时段更大的费用下降。

\subsubsection{波动电价影响的解释}

已有分析进一步将价格与净负荷的关系分解为日期尺度和日内尺度。在 365 天全量数据窗口下，日级价格水平与当日净负荷的相关系数为
0.9824；在本文 334 天评价窗口下，该相关系数为 0.9801。日内价格形状与日内净负荷形状的相关系数为
0.093。上述结果表明两种时间尺度具有不同的对应特征，日期尺度上的较强相关性说明，购电需求较大的日期也可能面临较高价格，从而增加评价期总支出。

储能通过跨时段转移电量缓解该影响，但其可用储电量范围为 9600
kWh，且受到最大充放电功率与转换损耗限制，无法任意转移全部购电需求。交叉结算结果显示，重新优化能够降低既定波动价格下的费用，而外部价格变化仍占据成本增加的主要部分。

因此，在本题数据与既定信息条件下，波动电价下的调度需要同时考虑价格与供需的时序关系。储能的作用体现为调整购电时段、利用光伏富余并缓解高成本时段的外购需求，其经济效果应通过完整费用核算评价。

\section{模型结果分析与检验}

\subsection{综合结果与一致性检验}

四问共用同一时间尺度、能量平衡和储能约束，并依次增加光伏不确定性、预报更新及波动电价。各费用差的解释对象不同：问题一比较同日有无储能，问题二比较同预测下是否设置因果裕度，问题三比较预报更新组合，问题四用交叉结算分解价格环境与策略重优化。因此，跨问结果只作机制衔接，不把不同信息集和交易规则下的费用差解释为单一算法收益。

物理层逐十分钟复算供需平衡和储电量递推，同时检查功率、容量与充放电互斥；问题一另检验日初日末相等，问题二至问题四检验跨日状态连续。四小时充放电量仅为24个十分钟时段的汇总，不直接与单时段上限比较。已有复算中，问题三最大平衡残差不超过 \(4.5\times10^{-13}\) kWh，2004个四小时汇总区间均与执行结果一致。

费用层分别核对计划量、最终调整量、实际调用量和紧急购电量。普通购电按各问交易规则结算，终端价值只参与决策、不计入现金收入。数据位置按144个时段依次对应0:00---24:00；紧急购电表中空白日期继承上一非空日期，避免漏读连续区间。

信息层通过改变当天实际光伏、保持计划时可用信息不变来检验非预见性：问题二和问题四的计划保持不变，而实际执行随光伏变化。将问题四价格恢复为附件1后，Q4-2与问题二费用仅有约0.002元数值差，Q4-3恢复为 \(13{,}120{,}194.06\) 元，与问题三一致。结合单元格扰动和整日错位测试，说明模型能够同时识别物理、时间映射与结果输出异常。
\subsection{灵敏度与稳健性分析}\label{sec:sensitivity}

为检验结论对关键参数与外部条件的依赖程度，分别改变终端价值、储能效率、预报误差强度和价格环境，每次只改变一个因素，其余信息、模型、求解与结算口径保持一致。评价期统一为2025年2月1日至12月31日共334天，各方案均从2月1日0:00的6000 kWh储电量出发，使差异可归因于被测因素。

评价统一采用两把尺度：实际现金费用 \(J_{\rm cash}\) 与固定终端价值 \(\lambda_{\rm eval}=0.4720\) 元/kWh 的库存修正费用
\begin{equation}
J_{\rm cmp}=J_{\rm cash}+\lambda_{\rm eval}\left(S_{\rm eval,start}-S_{\rm eval,end}\right).
\end{equation}
由于终端价值系数 \(\lambda\) 与储能效率 \(\eta\) 都会改变年末库存，仅看 \(J_{\rm cash}\) 会被``放空库存换现金''污染，因此 \(\lambda\) 与 \(\eta\) 两条轴以 \(J_{\rm cmp}\) 为主判据、\(J_{\rm cash}\) 为辅；预报误差轴 \(\rho\) 与价格机制轴的年末库存不变，直接以 \(J_{\rm cash}\) 比较。主要实验轴及测试水平见下表。

\begin{table}[!htbp]
\centering
\caption{灵敏度与稳健性分析的主要实验轴}
\footnotesize
\setlength{\tabcolsep}{2pt}
\begin{tabular}{|c|c|c|c|}
\hline
实验 & 测试轴 & 作用对象 & 测试水平 \\
\hline
S0 终端价值稳健性 & 终端价值系数 \(\lambda_{\rm dec}\) & Q2 & \(\{0,0.3776,0.4720,0.5664\}\) 元/kWh \\
\hline
S1 物理效率稳健性 & 单程效率 \(\eta_c=\eta_d=\eta\) & Q1/Q2/Q4-2（全年） & \(\{0.85,0.90,0.95\}\) \\
\hline
 &  & Q3/Q4-3（26 日筛查） & 同上 \\
\hline
S2 预报误差传导 & 预报误差缩放 \(\rho\) & Q3 & \(\{0.5,1.0,1.5\}\) \\
\hline
S3 价格机制分解 & 价格环境 \(\times\) 计划来源 & Q3/Q4-3 & 固定口径的2\(\times\)2交叉结算 \\
\hline
S4 既有消融证据 & 既定对照方案 & Q1--Q4 & 复核并统一呈现 \\
\hline
\end{tabular}
\end{table}

其中 S0--S3 为核心实验：S0 检验决策偏好参数，S1 检验共同物理参数，S2 为预报信息质量的反事实诊断，S3 为价格环境的机制消融；S4 复核既有对照。四者的扰动源互不重叠，合起来构成对``参数---信息---环境''完整外生面的覆盖。

\subsubsection{终端价值参数的敏感性}

问题二的日前目标中包含终端价值项，系数 \(\lambda_{\rm dec}\) 决定日末库存的局部价值。测试水平取
\begin{equation}
\lambda_{\rm dec}\in\{0,\ 0.3776,\ 0.4720,\ 0.5664\}\ {\rm 元/kWh},
\end{equation}
其中 0.4720 为基准，0.3776 与 0.5664 构成不使用正式期结果调参的对称 \(\pm20\%\) 局部扰动，0 用于检验忽略终端库存价值的极端情形。四档的全年结果见下表。

\begin{table}[!htbp]
\centering
\caption{问题二终端价值系数的敏感性结果（主判据 \(J_{\rm cmp}\)）}
\footnotesize
\begin{tabular}{|c|c|c|c|c|c|c|}
\hline
\(\lambda_{\rm dec}\) & \(J_{\rm cmp}\)（主） & \(\Delta J_{\rm cmp}\) & \% & \(J_{\rm cash}\)（辅） & \(\Delta J_{\rm cash}\) & 年末 SOC/kWh \\
\hline
0 & 13,253,490.73 & +2,070.04 & +0.016\% & 13,251,225.13 & $-$1,115.96 & 1200 \\
\hline
0.3776（$-$20\%） & 13,253,490.73 & +2,070.04 & +0.016\% & 13,251,225.13 & $-$1,115.96 & 1200 \\
\hline
\textbf{0.4720（基准）} & \textbf{13,251,420.69} & --- & --- & 13,252,341.09 & --- & \textbf{7950} \\
\hline
0.5664（+20\%） & 13,285,666.22 & +34,245.53 & +0.258\% & 13,287,931.82 & +35,590.73 & 10800 \\
\hline
\end{tabular}
\end{table}

结果呈“平台—跳变”结构：当 \(\lambda_{\rm dec}\leq0.3776\) 时，年末储电量降至1200 kWh，\(J_{\rm cmp}=13{,}253{,}490.73\) 元；基准系数0.4720对应7950 kWh；系数0.5664对应10800 kWh。在\(\pm20\%\)扰动内，\(J_{\rm cmp}\)最大变化34,245.53元，占基准的0.258\%。若只看现金费用，\(\lambda=0\)似乎节省1,115.96元，但其年末库存少6750 kWh；统一折算后反而增加2,070.04元。因此该参数应同时报告期末状态和\(J_{\rm cmp}\)，不宜用单一弹性概括阈值响应。
\subsubsection{储能效率的敏感性}

储能单程效率 \(\eta_c=\eta_d=\eta\) 同时改变可用能量与套利、应急之间的权衡，是四问共同的关键物理参数。本题基准为\textbf{单程} 0.90（对应往返 0.81），测试水平取
\begin{equation}
\eta\in\{0.85,\ 0.90,\ 0.95\},
\end{equation}
对应往返效率 \(\eta^{2}\in\{0.7225,\ 0.81,\ 0.9025\}\)；0.85 与 0.95 为围绕基准的工程压力水平，不表示设备效率的统计置信区间。效率在求解、实际调用、结果核验与输出等环节中同步改变，避免各环节口径不一致。

问题一在硬闭环 \(S_{0}=S_{144}=6000\) 条件下运行完整单日两阶段词典序 LP；问题二与问题四的日前计划分支运行完整评价期；问题三与问题四的滚动调整分支在 26 个代表性日期（2--12 月每月 1 日与 15 日加四个题目指定日期）上做固定日初状态的局部筛查，每个日期的日初 SOC 取自基线轨迹同日值。五类求解分支的结果见下表。

\begin{table}[!htbp]
\centering
\caption{储能单程效率敏感性结果（\(\lambda/\eta\) 轴主判据 \(J_{\rm cmp}\)）}
\footnotesize
\begin{tabular}{|c|c|c|c|c|c|}
\hline
问题 & 口径 & \(\eta=0.85\) & \(\eta=0.90\)（基准） & \(\eta=0.95\) & \(\Delta\%\)（0.85／0.95） \\
\hline
Q1 & 单日，\(J_{\rm actual}\) & 36,603.41 & 35,126.95 & 33,767.03 & +4.20\%／$-$3.87\% \\
\hline
Q2 & 全年，\(J_{\rm cmp}\) & 13,717,516.01 & 13,251,420.69 & 12,794,343.96 & +3.52\%／$-$3.45\% \\
\hline
Q4-2 & 全年，\(J_{\rm cmp}\) & 14,327,936.58 & 13,849,888.24 & 13,348,273.60 & +3.45\%／$-$3.62\% \\
\hline
Q3 & 26 日筛查，\(\Sigma J_{\rm cmp}\) & 1,006,573.29 & 973,053.04 & 936,051.09 & +3.44\%／$-$3.80\% \\
\hline
Q4-3 & 26 日筛查，\(\Sigma J_{\rm cmp}\) & 1,045,388.76 & 1,009,100.26 & 970,467.59 & +3.60\%／$-$3.83\% \\
\hline
\end{tabular}
\end{table}

单程效率由0.90降至0.85时，各分支费用上升3.4\%--4.2\%；升至0.95时下降3.4\%--3.9\%，方向一致。Q1因日末硬闭环无法跨日转移损失，变化幅度较大；Q3与Q4-3仅为26个代表日期的固定日初状态筛查，不能外推为全年幅度。效率变化还会改变年末库存，故同样以\(J_{\rm cmp}\)作为跨参数比较主判据。
\subsubsection{预报误差强度的决策传导}

附件 3 的预报误差强度通过离线反事实实验刻画。在发布时点、重构方式、滚动决策与执行规则全部不变的前提下，对原生小时节点按
\begin{equation}
F^{raw}_{d,r,T}(\rho)=P^{real}_{d,r,T}+\rho\left(F_{d,r,T}-P^{real}_{d,r,T}\right)
\end{equation}
构造误差缩放后的合成预测，其中 \(\rho\in\{0.5,1.0,1.5\}\) 分别表示误差减半、基准与放大 50\%。合成预测经过与基准Q3完全相同的左端点锚定、小时内插值、非负截断与时段末映射后进入滚动优化；离线生成与在线决策严格隔离，保证优化器只读取决策时刻已经发布的合成预测。该实验是使用实际出力构造误差尺度的离线反事实诊断，只刻画官方预报整体质量变化对调度成本、调整量与应急风险的影响，不作为可部署的预测方案。三档的全年结果见下表。

\begin{table}[!htbp]
\centering
\caption{问题三预报误差强度的决策传导（\(J_{\rm cash}\)，元）}
\footnotesize
\begin{tabular}{|c|c|c|c|c|c|c|}
\hline
\(\rho\) & \(J_{\rm cash}\) & \(\Delta\%\) & 计划基底 & 调整相关 & 紧急购电 & 紧急电量/kWh \\
\hline
0.5 & 12,805,045.44 & $-$2.40\% & 12,360,747.72 & 275,369.75 & 168,927.97 & 36,680.69 \\
\hline
\textbf{1.0（基准）} & \textbf{13,120,194.06} & --- & 12,381,239.55 & 542,052.30 & 196,902.21 & 42672.22 \\
\hline
1.5 & 13,478,837.67 & +2.73\% & 12,424,724.30 & 818,388.11 & 235,725.25 & 51,024.15 \\
\hline
\end{tabular}
\end{table}

预报误差强度减半或放大50\%时，全年费用分别变化\(-2.40\%\)与\(+2.73\%\)。计划基底仅变化\(-0.17\%\)与\(+0.35\%\)，调整相关成本却变化\(-49.2\%\)与\(+51.0\%\)，说明误差主要由调整层吸收；紧急购电费变化\(-14.2\%\)与\(+19.7\%\)，表明滚动调整缓冲了部分冲击。节点级MAE为93.30、186.61、279.91，严格按\(\rho\)缩放；槽级MAE为93.47、143.74、198.45 kW，非负截断槽数为130、76、112，从而同时核验扰动强度和时序重构。
\subsubsection{价格环境与策略重优化的机制分解}

前文已给出波动电价下两类策略的 \(2\times2\) 交叉结算，并以未舍入的全年结果完成主路径、反向路径与差上差交互的完整分解。沿用记法 \(J_{ab}\) 表示计划来源价格 \(a\)、结算价格 \(b\) 下的现金费用（1 表示附件 1 固定电价，4 表示附件 4 波动电价）。

% Q4交叉结算表的全部数据已合并至前文“问题四波动电价下的全年结果与归因”浮动体。

换成逐日逐时段波动电价后，问题三口径全年费用上升 606,839.60 元（+4.63\%）。其中 760,032.97 元来自价格环境本身（同一计划在新价格下结算的代价），重新优化抵消了 153,193.38 元，两者之比约为 5:1——价格环境的影响远大于重新优化的抵消能力。差上差交互项为 \(-258,149.70\) 元，说明价格环境效应\textbf{依赖计划来源}：同一份计划换价结算，与在新价下重算后再换价结算，两者的价差不可简单叠加。该交互项属于机制诊断，不解释为电价信息价值或算法改进。主路径与反向路径均闭合于总效应，分解残差为 0（未舍入值计算）。

\subsubsection{既有消融证据与稳健性边界}

前文问题二方案对照表表明：相对主方案，取消因果裕度使费用增加925,662.91元（+6.98\%）；反过来按无裕度方案为分母，主方案降费6.53\%。完全信息下界比主方案低997,644.54元（占主方案7.53\%），但不能在线实施；附件1典型日方案则高43.70\%。只有“同预测、裕度为0”的对照单独识别裕度作用，其余方案因预测方法或信息集改变，仅界定方法边界。

问题三由仅0:00更新扩展到四时点更新后，费用减少443,957.78元（3.27\%），紧急电量下降56.6\%；边际价值依次为12:00、6:00、18:00。问题四四种裕度结构的全年费用差异不超过\(\pm0.21\%\)，且Q4-2与Q2费用比值仅由1.0451变为1.0456，说明裕度结构不是跨问费用差的主要来源。
\begin{figure}[!htbp]
\centering
\includegraphics[width=0.95\textwidth]{q1_06_sensitivity.png}
\caption{问题一不同参数扰动下的购电费用变化}
\end{figure}
图15汇总问题一除效率三档外的22个既有扰动情景，覆盖负载、光伏、容量、功率与电价。由于各参数扰动幅度不同，图中仅判断同一参数内的方向和幅度，不作跨参数统一敏感度排序。

综合而言，终端价值主要形成库存阈值，效率变化在五类分支中方向一致，预报误差主要经调整层传导，价格环境效应约为策略重优化抵消额的5倍。以上均为同问题、同口径的单轴比较；跨问题费用差只用于机制衔接，不解释为单一算法收益。
\section{模型评价与改进}

\subsection{模型优点}

（1）\textbf{物理关系统一，能够衔接不同调度条件。}
四问采用一致的时间尺度、能量平衡和储能状态方程，在此基础上逐步引入光伏预测、购电调整及波动电价，便于解释不同条件下的策略变化。问题一采用两阶段词典序线性规划，在优先保证购电费用的基础上减少无效充放电，兼顾经济目标与运行合理性。

（2）\textbf{明确区分决策信息与实际执行。}
日前计划使用决策时刻可获得的光伏信息，实际执行根据真实出力与储电量处理偏差；滚动预报仅改变后续策略。计划量、调整量、实际调用量及紧急购电量分别定义，使信息条件、运行过程与费用核算相互对应。

（3）\textbf{结果分析能够解释费用变化的来源。}
通过储能状态价值分析说明边界电量的经济意义，通过预报更新对照评价新增信息的作用，并利用交叉结算区分价格环境变化与策略重优化的影响，使分析从总费用比较延伸至调度机制解释。

（4）\textbf{采用多层次检验支持结果可信性。}
已有验证覆盖能量守恒、储能边界、跨日连续性、信息可得性和输出映射，并通过独立复算、退化对照及人为错误测试检查计算与提交结果的一致性。

\subsection{模型局限与改进方向}

（1）\textbf{信息条件具有明确适用范围。}
模型假设当日负载已知，问题四主方案还假设当日电价已知。若实际应用中负载和价格同样需要预测，则需引入相应误差描述，并在一致的信息条件下重新评价策略。

（2）\textbf{预测与参数标定仍受历史样本限制。}
问题二采用短期历史均值与保守裕度，便于实施和复现，但对突发天气变化的适应能力有限。问题三的参数化结果及预报更新收益也与给定样本相关，不能直接推广为其他运行环境下的最优配置。后续可利用新增数据进行独立时段验证，考察参数与预测规则的稳定性。

（3）\textbf{储能成本描述有所简化。}
模型考虑充放电效率、容量和功率边界，未纳入自放电、容量衰减及循环寿命成本。问题一减少充放电吞吐量能够抑制无效循环，但不能替代完整的寿命成本模型。若获得相应参数，可进一步评价短期购电节省与长期设备损耗之间的关系。

（4）\textbf{终端价值与结算解释影响结果适用性。}
常数终端价值是对剩余电量后续作用的简化，不同终端设置可能形成不同期末库存。滚动调整采用既定结算口径，其经济结果亦依赖该解释。后续比较应保持边界状态和结算规则一致，并结合已有敏感性结果报告差异。

上述改进方向均属于模型的适用性扩展，不计入本文已完成方案的计算结果。

\section{AI 工具使用声明}

本参赛队在竞赛过程中使用了 AI 工具，主要用于建模思路论证、算法与代码实现及调试、结果审查与论文文字整理，详细使用情况见支撑材料。

\clearpage
{\small
\begin{thebibliography}{99}

\bibitem{arrow1951} ARROW K J, HARRIS T, MARSHAK J. Optimal inventory policy[J]. Econometrica, 1951, 19(3): 250-272.

\bibitem{bellman1957} BELLMAN R. Dynamic Programming[M]. Princeton: Princeton University Press, 1957.

\bibitem{bental2004} BEN-TAL A, GORYASHKO A, GUSLITZER E, et al. Adjustable robust solutions of uncertain linear programs[J]. Mathematical Programming, 2004, 99(2): 351-376.

\bibitem{birge2011} BIRGE J R, LOUVEAUX F. Introduction to Stochastic Programming[M]. 2nd ed. New York: Springer, 2011.

\bibitem{celades2024} CELADES E, PÉREZ E, APARICIO N, et al. Tool for optimization of sale and storage of energy in wind farms[J]. Mathematics and Computers in Simulation, 2024, 224: 2-18. DOI: 10.1016/j.matcom.2023.03.010.

\bibitem{chen2011} CHEN C, DUAN S, CAI T, et al. Smart energy management system for optimal microgrid economic operation[J]. IET Renewable Power Generation, 2011, 5(3): 258-267. DOI: 10.1049/iet-rpg.2010.0052.

\bibitem{conejo2010} CONEJO A J, CARRIÓN M, MORALES J M. Decision Making Under Uncertainty in Electricity Markets[M]. New York: Springer, 2010.

\bibitem{daneshvar2020} DANESHVAR M, MOHAMMADI-IVATLOO B, ZARE K, et al. Two-stage robust stochastic model scheduling for transactive energy based renewable microgrids[J]. IEEE Transactions on Industrial Informatics, 2020, 16(11): 6857-6867. DOI: 10.1109/TII.2020.2973740.

\bibitem{dantzig1963} DANTZIG G B. Linear Programming and Extensions[M]. Princeton: Princeton University Press, 1963.

\bibitem{dufolpez2015} DUFO-LÓPEZ R. Optimisation of size and control of grid-connected storage under real time electricity pricing conditions[J]. Applied Energy, 2015, 140: 395-408. DOI: 10.1016/j.apenergy.2014.12.012.

\bibitem{fazlalipour2019} FAZLALIPOUR P, EHSAN M, MOHAMMADI-IVATLOO B. Risk-aware stochastic bidding strategy of renewable micro-grids in day-ahead and real-time markets[J]. Energy, 2019, 171: 689-700. DOI: 10.1016/j.energy.2018.12.173.

\bibitem{grimaldi2024} GRIMALDI A, MINUTO F D, BROUWER J, et al. Profitability of energy arbitrage net profit for grid-scale battery energy storage considering dynamic efficiency and degradation using a linear, mixed-integer linear, and mixed-integer non-linear optimization approach[J]. Journal of Energy Storage, 2024, 95: 112380. DOI: 10.1016/j.est.2024.112380.

\bibitem{gulin2017} GULIN M, PAVLOVIĆ T, VAŠAK M. A one-day-ahead photovoltaic array power production prediction with combined static and dynamic on-line correction[J]. Solar Energy, 2017, 142: 49-60. DOI: 10.1016/j.solener.2016.12.008.

\bibitem{howard1966} HOWARD R A. Information value theory[J]. IEEE Transactions on Systems Science and Cybernetics, 1966, 2(1): 22-26.

\bibitem{hyndman2018} HYNDMAN R J, ATHANASOPOULOS G. Forecasting: Principles and Practice[M]. 2nd ed. Melbourne: OTexts, 2018.

\bibitem{kaufman1990} KAUFMAN L, ROUSSEEUW P J. Finding Groups in Data: An Introduction to Cluster Analysis[M]. New York: Wiley, 1990.

\bibitem{kleywegt2002} KLEYWEGT A J, SHAPIRO A, HOMEM-DE-MELLO T. The sample average approximation method for stochastic discrete optimization[J]. SIAM Journal on Optimization, 2002, 12(2): 479-502.

\bibitem{li2022} LI Y, PENG J, JIA H, et al. Optimal battery schedule for grid-connected photovoltaic-battery systems of office buildings based on a dynamic programming algorithm[J]. Journal of Energy Storage, 2022, 50: 104557. DOI: 10.1016/j.est.2022.104557.

\bibitem{lujanorojas2017} LUJANO-ROJAS J M, DUFO-LÓPEZ R, BERNAL-AGUSTÍN J L, et al. Optimizing daily operation of battery energy storage systems under real-time pricing schemes[J]. IEEE Transactions on Smart Grid, 2017, 8(1): 316-330. DOI: 10.1109/TSG.2016.2602268.

\bibitem{malysz2013} MALYSZ P, SIROUSPOUR S, EMADI A. MILP-based rolling horizon control for microgrids with battery storage[C]//IECON 2013 - 39th Annual Conference of the IEEE Industrial Electronics Society. Vienna: IEEE, 2013: 2099-2104.

\bibitem{nadaraya1964} NADARAYA E A. On estimating regression[J]. Theory of Probability and Its Applications, 1964, 9(1): 141-142.

\bibitem{parisio2014} PARISIO A, RIKOS E, GLIELMO L. A model predictive control approach to microgrid operation optimization[J]. IEEE Transactions on Control Systems Technology, 2014, 22(5): 1813-1827. DOI: 10.1109/TCST.2013.2295737.

\bibitem{sioshansi2009} SIOSHANSI R, DENHOLM P, JENKIN T, et al. Estimating the value of electricity storage in PJM: arbitrage and some welfare effects[J]. Energy Economics, 2009, 31(2): 269-277. DOI: 10.1016/j.eneco.2008.10.005.

\bibitem{umezozor2016} UMEZOZOR E C, TRIFKOVIC M. Operational scheduling of microgrids via parametric programming[J]. Applied Energy, 2016, 180: 672-681. DOI: 10.1016/j.apenergy.2016.08.009.

\bibitem{watson1964} WATSON G S. Smooth regression analysis[J]. Sankhyā: The Indian Journal of Statistics, Series A, 1964, 26(4): 359-372.

\bibitem{wu2018} WU J, XING X, LIU X, et al. Energy management strategy for grid-tied microgrids considering the energy storage efficiency[J]. IEEE Transactions on Industrial Electronics, 2018, 65(12): 9539-9549. DOI: 10.1109/TIE.2018.2818660.

\bibitem{xu2010} XU Y, XIE L, SINGH C. Optimal scheduling and operation of load aggregator with electric energy storage in power markets[C]//2010 North American Power Symposium. Arlington, TX: IEEE, 2010: 1-7.

\bibitem{zhang2023} ZHANG X, HUANG C, SHEN J. Energy optimal management of microgrid with high photovoltaic penetration[J]. IEEE Transactions on Industry Applications, 2023, 59(1): 128-137. DOI: 10.1109/TIA.2022.3208885.

\bibitem{zhang2018} ZHANG Y, MENG F, WANG R, et al. A stochastic MPC based approach to integrated energy management in microgrids[J]. Sustainable Cities and Society, 2018, 41: 349-362. DOI: 10.1016/j.scs.2018.05.044.

\end{thebibliography}
}
\end{document}
